\unnumberedchapter{Résumé de la thèse}

\section*{Introduction}
Ma thèse porte sur l'étude du mouvement brownien. En tant que scientifiques et a fortiori physiciens, nous nous familiarisons avec ce phénomène dès nos années lycéennes, tant il est parfaitement connu et prévisible pour l'étude de cas très simples. On nous fournit, typiquement, une solution aqueuse dans laquelle se trouvent de très petites particules (des colloïdes). Nous les observons au microscope optique et nous nous étonnons, parce que notre culture scientifique n'est pas encore tout à fait développée, du mouvement inarrêtable et imprévisible des colloïdes, alors que le fluide les baignant est tout à fait quiescent. Nous apprenons ensuite que ce que nous observons n'est autre qu'une des premières preuves expérimentales de la nature atomique de la matière. Notre instrument est très limité dans ses capacités d'observation, et sa limite basse en termes de taille se situe plusieurs ordres de grandeur au-delà de la taille d'une molécule. Nous voyons, en fait, la conséquence de l'existence de ces atomes, plutôt que les atomes eux-mêmes : c'est la température du fluide elle-même qui se traduit en mouvement des molécules à leur échelle. Ces molécules, bien que minuscules devant la taille d'un colloïde, sont animées d'une vitesse élevée, de plusieurs centaines de mètres par seconde en moyenne. En heurtant le colloïde, elles lui transfèrent à chaque impact un peu de leur quantité de mouvement, et comme ces impacts eux-mêmes sont aléatoires tant dans le temps que dans la direction, le colloïde s'anime d'un mouvement lui-même imprévisible et erratique.

Ces observations ont été largement confirmées et formalisées par d'éminents chercheurs, et je décris le cas historique du mouvement brownien simple dans le premier chapitre de ce manuscrit. Je précise aussi dans ce chapitre les différents développements théoriques importants pour notre bonne compréhension de mon travail de thèse et y fais référence tout au long du manuscrit. En outre, le mouvement brownien n'est pas qu'un simple cas d'école, fabriqué pour nous faire déduire la réalité moléculaire. Il s'agit, en fait, d'un processus extrêmement commun dans de nombreux systèmes biologiques. Si nous reprenons sa première observation historique dans des grains de pollen, nous notons déjà que la forme des colloïdes elle-même entraîne des modifications aux théories canoniques; cela constitue d'ailleurs une des premières généralisations apportées à la théorie du mouvement brownien. L'on peut aussi se questionner, puisque ce mouvement naît de l'agitation thermique au sein du fluide, de ce qu'il en advient quand le fluide lui-même est perturbé, par exemple en modifiant les propriétés mêmes (viscosité, température), ou encore en quittant la limite du cas très dilué où les colloïdes sont traités comme infiniment éloignés, dans un régime dense où le fluide déplacé par le mouvement de l'un affecte directement celui de l'autre et où un couplage hydrodynamique apparaît. Une autre exploration sur l'effet de parois rigides, là encore modifiant le comportement du fluide et donc du mouvement du colloïde, vient tout aussi naturellement. Ces questions ont été explorées et en partie résolues dans les cent dernières années, mais leur complexité ne permet pas toujours de comprendre totalement et parfaitement leur unification. J'explore cette question dans le chapitre 2. Puis vient la question de l'influence de la déformabilité dans des systèmes confinés. Le vivant regorge de membranes déformables telles qu'une paroi cellulaire, une muqueuse ou un tissu mou. Parfois, ces membranes elles-mêmes sont sujettes à ce même mouvement provenant des chocs avec les molécules de fluide. Comment prédire l'influence de leur déformabilité et de leurs fluctuations sur un colloïde voisin ? Peut-on contrôler, améliorer ou diminuer la mobilité d'une molécule à des fins thérapeutiques ? Au contraire, serait-il possible de stopper la propagation d'un pathogène sur de telles membranes, responsable de la dégénérescence parfois incurable de certaines de nos capacités ? Ce problème est par essence très complexe à résoudre, mais je décris dans le chapitre 4 un modèle élémentaire et théorique pour capturer qualitativement l'effet de la déformabilité sur le mouvement brownien. L'écoulement du fluide lui-même affecte lui aussi le mouvement des colloïdes, d'abord et évidemment puisque ce dernier se retrouve entraîné dans l'écoulement, mais aussi pour des raisons moins triviales. En outre, je cherche à comprendre la possibilité ou non d'un effet de la densité, même faible, d'une solution de colloïdes dans un écoulement sur la dynamique de ceux-ci. Enfin, nous abordons déjà certes de nombreux cas de figure propres au monde du vivant, mais \emph{quid} du vivant lui-même ? Les colloïdes eux-mêmes peuvent dans certains cas se déplacer par leurs propres moyens en utilisant l'énergie disponible qui les entoure. S'ils sont parfois trop gros pour être visiblement affectés par l'agitation thermique du fluide, nous pouvons alors prédire leur comportement en les traitant par extension comme des particules browniennes effectives. Il est alors possible d'extraire encore de l'information de leurs trajectoires, tout comme nous le ferions pour un colloïde passif de petite taille. J'utilise dans le chapitre 5 ce type de modèle pour étudier le mouvement de micronageurs d'une taille bien supérieure à celle d'une particule brownienne, afin notamment d'étudier les comportements de micronageurs dans un écoulement comme décrit dans le chapitre 4.


\section*{Libre accès}
L'ensemble des résultats que j'ai produits tout au long de ma thèse s'appuie sur la simulation numérique comme principal outil. J'ai choisi au début de ma thèse et \emph{ex nihilo} d'aborder la recherche par cet angle, car il est le seul permettant de proprement isoler différents effets pour mieux prédire leurs conséquences. Il ne serait pas possible, expérimentalement, d'observer une sphère parfaite, un zéro absolu, un écoulement strictement parabolique ou une eau infiniment pure. Pour moi, il ne s'agit à chaque fois que de quelques lignes de code tout au plus. Naturellement, avec cette capacité vient la question de la réalité de ce que je simule. Je dois toujours garder un œil extrêmement sceptique sur mes résultats, et les comparer aux théories et expériences, quand celles-ci existent. Je vois l'outil numérique comme un fantastique couteau suisse, permettant à la fois l'exploration d'effets inaccessibles expérimentalement à cause des limites actuelles de la technologie, et la validation de théories irréalisables à cause des lois de la nature, alors qu'elles apportent tout de même, sinon quantitativement, au moins conceptuellement, des développements vitaux pour la recherche.

Puisque je me suis formée tout au long de ma thèse grâce aux ressources (forums, publications ou tutoriels) gratuites en libre accès, ainsi qu'aux interactions humaines qui sont deux piliers de la recherche, tous les codes que j'ai développés sont disponibles publiquement sur \href{https://github.com/Laacherez/thesis}{une page GitHub}\footnote{https://github.com/Laacherez/thesis/Codes}. J'accorde une immense importance à la science libre, et j'estime que la réussite d'un chercheur ne doit pas se limiter aux fonds qu'on lui accorde, mais à sa curiosité, à sa capacité à poser les bonnes questions, à s'informer et à collaborer pour les résoudre, et à sa pédagogie et son intégrité dans la communication de ses résultats.